% 三视图横向并排:正视图/侧视图(左)/俯视图;线宽.6pt,12pt
% 点修: 主视图->正视图,左视图->侧视图(左);几何/尺寸/线型未动,全实线
\begin{tikzpicture}[line width=.6pt,font=\normalsize,x=1cm,y=1cm]
% 正视图:矩形宽√2高1(绘图×2:2.828×2)
\begin{scope}[shift={(0,3.2)}]
\draw (0,0) rectangle (2.828,2);
\draw[thin] (0,2) -- (0,2.5);
\draw[thin] (2.828,2) -- (2.828,2.5);
\draw[<->] (0,2.3) -- node[above,fill=white,inner sep=1pt] {$\sqrt{2}$} (2.828,2.3);
\draw[thin] (2.828,0) -- (3.328,0);
\draw[thin] (2.828,2) -- (3.328,2);
\draw[<->] (3.128,0) -- node[right,fill=white,inner sep=1pt] {$1$} (3.128,2);
\node[below] at (1.414,-0.3) {正视图};
\end{scope}
% 侧视图(左):直角三角形,竖直边(左)高2,底边宽2.828,斜边左上至右下
\begin{scope}[shift={(4.6,3.2)}]
\draw (0,0) -- (0,2) -- (2.828,0) -- cycle;
\node[below] at (1.414,-0.3) {侧视图（左）};
\end{scope}
% 俯视图:正方形边长2.828
\begin{scope}[shift={(9.0,3.2)}]
\draw (0,0) rectangle (2.828,2.828);
\draw[thin] (2.828,0) -- (3.328,0);
\draw[thin] (2.828,2.828) -- (3.328,2.828);
\draw[<->] (3.128,0) -- node[right,fill=white,inner sep=1pt] {$\sqrt{2}$} (3.128,2.828);
\node[below] at (1.414,-0.3) {俯视图};
\end{scope}
\end{tikzpicture}
